\documentclass[twoside,openany]{book}
\usepackage[czech]{babel}			% Language
\usepackage[utf8]{inputenc}			% Encoding of characters in this .tex file
\usepackage{cmap}					% Make PDF file searchable and copyable (ASCII characters)
\usepackage{lmodern}				% Make PDF file searchable and copyable (Accented characters)
\usepackage[T1]{fontenc}			% Hyphenate accented words
\usepackage[a6paper, top=0.75cm, bottom=1.5cm, inner=0.75cm, outer=0.75cm]{geometry}		% Paper size and margins
\usepackage[fontsize=9]{scrextend}  % Font size
\usepackage{marginnote}				% Enable margin notes
\usepackage{rotating}				% Rotated environments (for margin headers)
\usepackage[protrusion]{microtype}	% Better typeset results
\usepackage{enumitem}				% Nicer enumeration lists
\usepackage{ragged2e}				% Add \justify command
\usepackage{titlesec}				% Customize chapters
\usepackage[usenames,dvipsnames,svgnames,table]{xcolor}		% Custom colors
\usepackage{tikzsymbols}			% Smileys
\usepackage{chronology}				% Timeline
\usepackage{tikz}                   % Number lines
%\usepackage{fontawesome}			% Special symbols

% ======= CUSTOM DOCUMENT SETTINGS ========

% Specify relevant Tikz libraries
\usetikzlibrary{arrows}

% Change \chapter and \section command display and spacings
% Source: https://tex.stackexchange.com/questions/63390/how-to-decrease-spacing-before-chapter-title
\titleformat{\chapter}[display]{\normalfont\Large\bfseries}{Týden \thechapter}{0em}{\Large}
\titlespacing*{\chapter}{0em}{-2em}{0em}
\titleformat{\section}[display]{\normalfont\bfseries}{\sectionname}{}{}
\titlespacing*{\section}{0em}{0.5em}{0em}

% \punct command allows to shift the footnotemark above a punctuation
\newlength{\spc} % declare a variable to save spacing value
\newcommand{\punct}[1]{%
  \settowidth{\spc}{#1}% set value of \spc variable to the width of #1 argument
  \addtolength{\spc}{-1.8\spc}% subtract from \spc about two (1.8) of its values making its magnitude negative
  #1% print the optional argument
  \hspace*{\spc}% print an additional negative spacing stored in \spc after #1
}

% Remove headers and position page number in the bottom mid
\pagestyle{plain}

% Smaller spaces between enumeration lists
\setitemize{itemsep=0em}

% Disable automatic indentation, enable paragraph skip
\setlength{\parindent}{0em}
\setlength{\parskip}{0.5em}

% Rule length including the margin
\newlength{\rulelength}
\setlength{\rulelength}{\dimexpr(\textwidth+0.5cm)\relax}

% New table column types with flexible width
%\newcolumntype{C}{>{\centering\arraybackslash}X}
\newcolumntype{C}[1]{>{\centering\arraybackslash}m{#1}}

% Rotated table cell
\newcommand{\rotatedCell}[2][l]{\rotatebox{90}{\renewcommand{\arraystretch}{0.75}\begin{tabular}[#1]{@{}l}#2\end{tabular}}}

% ======= CUSTOM ENVIRONMENTS ========

% Mala poznamka
\newcommand{\note}[1]{\textcolor{gray}{\small\itshape #1}}

% Poznamka za okrajom
\newcommand{\marginsection}[3]{\marginpar{\raisebox{#2\height}{
\begin{turn}{#1}
%\fontfamily{phv}\selectfont
\bfseries \color{gray} \large
#3\end{turn}}}}

% Casova os pre rozvrhnutie planu cvika
\newcommand{\timeline}{
\vspace*{2.1cm}
\color{gray}
\begin{chronology}[15]{0}{120}{\textwidth}
\end{chronology}
\vspace*{-0.2cm}
\color{black}
}

% Smajlici pre sebahodnotenie pocitu z cvika
\newcommand{\smileys}{
Jak jsem spokojen(a) s hodinou?
\raisebox{-0.3em}{\,\,\Changey[2]{-2}\,\,\Changey[2]{-1}\,\,\Changey[2]{0}\,\,\Changey[2]{1}\,\,\Changey[2]{2}}
}

% Zhodnotenie 2 kladov a 2 zaporov cvika
\newcommand{\goodbadpoints}{
Co se mi povedlo?
\begin{enumerate}
\item
\item
\end{enumerate}

Co mohlo být lépe?
\begin{enumerate}
\item
\item
\end{enumerate}
}

% Klasicky tyzden
\newcommand{\commonweek}{
\chapter{}
\vspace*{-2em}

\marginsection{90}{-1}{před výukou}

Příprava:\hspace{1cm} hodin(y)

Jaký je cíl této hodiny? Proč chci studenty učit tohle?
\vspace*{1cm}

Jak poznám, že se cíle podařilo dosáhnout?

\hspace*{-1cm}
\rule{\rulelength}{0.4pt}

\marginsection{90}{-3}{po výuce}

\smileys

\goodbadpoints

Vlastní otázky:

\newpage
\marginsection{-90}{1}{volné poznámky}
}

% Stranka rubriky
\newcommand{\rubricpage}[8]{
\newpage
\stepcounter{rubricquestion}
\section*{\therubricquestion: #1}

\smallskip
\note{Neznalý:} #2

\note{Začátečník:} #3

\note{Guru:} #4

\bigskip
\begin{tikzpicture}
\draw (0,0) -- (8,0);
\foreach \i in {0,1,...,8} % numbers on line
{
\fill[black] (\i,0) circle (1.5 mm);
\fill[white] (\i,0) circle (1.4 mm);
}
\node at (0.15, -0.5) {neznalý};
\node at (3, -0.5)    {začátečník};
\node at (8, -0.5)    {guru};
\end{tikzpicture}

\rule{\textwidth}{0.4pt}

\stepcounter{rubricquestion}
\section*{\therubricquestion: #5}

\smallskip
\note{Neznalý:} #6

\note{Začátečník:} #7

\note{Guru:} #8

\bigskip
\begin{tikzpicture}
\draw (0,0) -- (8,0);
\foreach \i in {0,1,...,8} % numbers on line
{
\fill[black] (\i,0) circle (1.5 mm);
\fill[white] (\i,0) circle (1.4 mm);
}
\node at (0.15, -0.5) {neznalý};
\node at (3, -0.5)    {začátečník};
\node at (8, -0.5)    {guru};
\end{tikzpicture}
}

% ======= TITLE PAGE =======

\title{Reflektivní deník učitele}
\author{Kolektiv vyučujících Praktika vedení cvičení}
\date{\today}

% ======= DOCUMENT START =======

\begin{document}
\begin{titlepage}
	\centering
    \vspace*{1cm}
	\includegraphics[width=0.5\textwidth]{img/diary}\par
	\vspace{0.5cm}
	{\huge\bfseries Reflektivní\\ deník učitele\par}
	\vfill
	\raggedright
    \Large
	Majitel:\par
    \vspace{0.2cm}
    Předmět:\par
\end{titlepage}

% ======= INTRODUCTION =======

\chapter*{Jak se z člověka stane dobrý učitel?}

\vspace*{1em}
Položme si namísto statické otázky dynamickou:

\vspace*{1em}
\textit{\large \uv{Co můžeš dělat pro to, aby ses jako učitel zlepšoval?}}

\vspace*{1em}
Ukazuje se, že jsou nutné, a nejspíš i dostačující, tři věci\punct{.}\footnotemark
\footnotetext{Jsou i další věci, které mohou pomoci (například expertní zpětná vazba), ale ty nejsou dobře škálovatelné.}
\begin{enumerate}
\item \textbf{Pravidelně učit.} \\Potřebuješ mít vlastní zkušenost, nejlépe často a pravidelně opakovanou.
\item \textbf{Reflektovat svoje učení.} \\Všímej si a přemýšlej nad tím, co fungovalo dobře a co příště zkusíš jinak.
\item \textbf{Pozorovat, jak učí ostatní.} \\Přemýšlej o tom, jak to dělají, co jim funguje, co ne a co z toho můžeš použít.
\end{enumerate}

Tento deníček ti pomůže reflektovat způsob jakým učíš. Také v~něm nalezneš informace o tom, čeho je užitečné si všímat, jaké otázky si klást a nad čím přemýšlet. Není to však učebnice dobrého učitele -- je to pouze průvodce na cestě, po které už jdeš.

Vhodným doplňkem k vlastní reflexi výuky jsou diskuze s ostatními učiteli o tom, jak učí. Pokud se chceš setkávat s dalšími cvičícími pravidelně a navíc se dozvědět více o tom, co funguje a co ne, zapiš si předmět \textit{DUCIT Praktikum vedení cvičení} na Fakultě informatiky MU.

\newpage
\section*{Jak reflektivní deníček používat?}

Nejlépe pravidelně.

Při přípravě i po každé své odvedené výuce si zapiš svoje myšlenky. Na třináct týdnů v semestru máš nachystaných třináct dvoustránek, a na každé pár nabídnutých otázek. Deníček je malý mimo jiné i proto, aby ti psaní nezabralo moc času.

Za stránkami pro jednotlivé týdny najdeš ještě rubriku učitelských schopností, seznam indikátorů, které je možné na hodině sledovat, připomenutí některých užitečných nástrojů k výuce a na závěr prostor na své vlastní poznámky.

Někde záměrně používáme anglické termíny (jako například \textit{self-assessment}), protože vhodný český překlad buď neexistuje, nebo má konotace, které nechceme.

\section*{Jaké jsou cíle používání deníku?}

Chtěli bychom, aby ti práce s deníkem:
\begin{enumerate}[topsep=0pt]
\item Připomínala zamyslet se nad svou hodinou (reflexe).
\item Dávala prostor pro snadné psaní si poznámek pro příště.
\item Pomáhala uvědomit si, co všechno patří k tvojí roli učitele.
\item Umožňovala sledovat svůj vlastní pokrok.
\end{enumerate}

\section*{Máš nápady na zlepšení?}

Budeme rádi za jakékoliv nápady, zkušenosti či poznámky k deníku. Ozvi se prosím emailem Martinovi na \textit{mukrop@mail.muni.cz}.

\newgeometry{top=0.75cm, bottom=1.5cm, inner=0.75cm, outer=1.2cm,
			heightrounded, marginparwidth=0.5cm, marginparsep=0.3cm}

% ======= WEEK 1 =======

\chapter{}
\vspace*{-2em}

\marginsection{90}{-1.3}{vyplnit před výukou}

Příprava:\hspace{1cm} hodin(y)\\
\note{Příprava materiálů, úkolů, pokynů, ...}

Z jakých bloků se moje hodina skládá?\\
\note{Rozvrhni si 2--6 bloků na časovou osu níže.}

\timeline

Jaké precedenty chci ve výuce nastavit? \\
\note{Oslovování, (ne)formálnost, kladení otázek, začínání včas, ...}
\vspace*{1cm}

\hspace*{-1cm}
\rule{\rulelength}{0.4pt}

\marginsection{90}{-1.5}{vyplnit po výuce}

\smileys
\note{Jaký mám pocit já? Co vidím na studentech?}

\goodbadpoints

\newpage
\marginsection{-90}{1}{volné poznámky}
\vspace*{-2em}
\note{Jaké chtěné/nechtěné precedenty se objevily na hodině?\\
Jak na další hodině udržím ty chtěné a napravím ty nechtěné?\\
Jaká byla atmosféra během vyučování?\\
Rozumí studenti organizaci kurzu a ví, co se od nich očekává?\\
(Další tipy na otázky najdeš na straně \pageref{indikatory}.)}

% ======= WEEK 2 =======

\chapter{}
\vspace*{-2em}

\marginsection{90}{-1.2}{vyplnit před výukou}

Příprava:\hspace{1cm} hodin(y)\\
\note{Příprava materiálů, úkolů, pokynů, ...}

Z jakých bloků se moje hodina skládá?\\
\note{Rozvrhni si 2--6 bloků na časovou osu níže.}

\timeline

Jaké precedenty chci ve výuce udržet? \\
\note{Oslovování, (ne)formálnost, kladení otázek, začátek včas, ...}
\vspace*{1cm}

\hspace*{-1cm}
\rule{\rulelength}{0.4pt}

\marginsection{90}{-1.3}{vyplnit po výuce}

\smileys
\note{Jaký mám pocit já? Co vidím na studentech?}

\goodbadpoints

\newpage
\marginsection{-90}{1}{volné poznámky}
\vspace*{-2em}
\note{Na jaký aspekt výuky se zaměřím při nejbližší hospitaci?\\
Jaký aspekt výuky má sledovat kolega, který navštíví mou hodinu?\\
(Další tipy na otázky najdeš na straně \pageref{indikatory}.)}

% ======= WEEK 3 =======

\chapter{}
\vspace*{-2em}

\marginsection{90}{-1.1}{vyplnit před výukou}

Příprava:\hspace{1cm} hodin(y)

Z jakých bloků se moje hodina skládá?

\timeline

Jaké otázky na hodině položím skupině?
\vspace*{0.8cm}

\hspace*{-1cm}
\rule{\rulelength}{0.4pt}

\marginsection{90}{-1.5}{vyplnit po výuce}

\smileys

\goodbadpoints

Na které otázky skupina reagovala? Na které ne?

\newpage
\marginsection{-90}{1}{volné poznámky}
\vspace*{-2em}
\note{Kolik otázek do publika jsem položil(a)?\\
Kolik otázek položili studenti mně? Je to málo/dost/příliš?\\
(Další tipy na otázky najdeš na straně \pageref{indikatory}.)}

% ======= WEEK 4 =======

\chapter{}
\vspace*{-2em}

\marginsection{90}{-1}{vyplnit před výukou}

Příprava:\hspace{1cm} hodin(y)

Z jakých bloků se moje hodina skládá?

\timeline

\hspace*{-1cm}
\rule{\rulelength}{0.4pt}

\marginsection{90}{-1.6}{vyplnit po výuce}

\smileys

\goodbadpoints

\note{Sem napiš otázku, na kterou si chceš po výuce odpovědět.}\\
Vlastní otázka:

\newpage
\marginsection{-90}{1}{volné poznámky}
\vspace*{-2em}
\note{Jaké znalosti a dovednosti chci naučit? Popiš je ve formátu rubriky, tedy pojmenuj dovednost a popiš škálu (viz str.\ \pageref{rubrika}).\\
Jaký typ učitele chci být? Co bych měl(a) pro to dělat?\\
(Další tipy na otázky najdeš na straně \pageref{indikatory}.)}

% ======= WEEK 5 =======

\chapter{}
\vspace*{-2em}

\marginsection{90}{-1.2}{před výukou}
Příprava:\hspace{1cm} hodin(y)

Jaký je cíl této hodiny? Proč chci studenty učit tohle? \\
\note{Při formulování cílů ti pomůžou slovesa na straně \pageref{bloom}.}
\vspace*{1cm}

Jak poznám, že se cíle podařilo dosáhnout? \\
\note{Existuje test, kterým je možné zjistit, zda jsi cíle dosáhl(a)?}

\hspace*{-1cm}
\rule{\rulelength}{0.4pt}

\marginsection{90}{-2.5}{po výuce}

\smileys

\goodbadpoints

\note{Sem napiš otázku, na kterou si chceš po výuce odpovědět.}\\
Vlastní otázka:

\newpage
\marginsection{-90}{1}{volné poznámky}
\vspace*{-2em}
\note{Podařilo se mi naplnit cíl hodiny?\\
Co je třeba do budoucna na téhle hodině změnit?\\
Jak budu plánovat a strukturovat nejbližší hodinu?\\
(Další tipy na otázky najdeš na straně \pageref{indikatory}.)}

% ======= WEEK 6 =======

\chapter{}
\vspace*{-2em}

\marginsection{90}{-1.2}{před výukou}

Příprava:\hspace{1cm} hodin(y)

Jaký je cíl této hodiny? Proč chci studenty učit tohle? \\
\note{Při formulování cílů ti pomůžou slovesa na straně \pageref{bloom}.}
\vspace*{1cm}

Jak poznám, že se cíle podařilo dosáhnout? \\
\note{Existuje test, kterým je možné zjistit, zda jsi cíle dosáhl(a)?}

\hspace*{-1cm}
\rule{\rulelength}{0.4pt}

\marginsection{90}{-2.5}{po výuce}

\smileys

\goodbadpoints

\note{Sem napiš otázku, na kterou si chceš po výuce odpovědět.}\\
Vlastní otázka:

\newpage
\marginsection{-90}{1}{volné poznámky}
\vspace*{-2em}
\note{Podařilo se mi naplnit cíl hodiny?\\
Co je třeba do budoucna na téhle hodině změnit?\\
Co je problematické na úkolech, které jsem zadal?\\
(Další tipy na otázky najdeš na straně \pageref{indikatory}.)}

% ======= WEEK 7 =======
\commonweek

% ======= WEEK 8 =======
\commonweek

% ======= WEEK 9 =======
\commonweek

% ======= WEEK 10 =======
\commonweek

% ======= WEEK 11 =======
\chapter{}
\vspace*{-2em}

\marginsection{90}{-1}{před výukou}

Příprava:\hspace{1cm} hodin(y)

Jaký je cíl této hodiny? Proč chci studenty učit tohle?
\vspace*{1cm}

Jak poznám, že se cíle podařilo dosáhnout?

\hspace*{-1cm}
\rule{\rulelength}{0.4pt}

\marginsection{90}{-3}{po výuce}

\smileys

\goodbadpoints

Blíží se konec semestru -- dostávají studenti to, co potřebují pro praxi/ke zkoušce?\\
\note{Jestli ne, jaké změny můžu v následujících týdnech provést?}

\newpage
\marginsection{-90}{1}{volné poznámky}

% ======= WEEK 12 =======

\chapter{}
\vspace*{-2em}

\marginsection{90}{-1}{před výukou}

Příprava:\hspace{1cm} hodin(y)

Jaký je cíl této hodiny? Proč chci studenty učit tohle?
\vspace*{1cm}

Jak poznám, že se cíle podařilo dosáhnout?

\hspace*{-1cm}
\rule{\rulelength}{0.4pt}

\marginsection{90}{-3}{po výuce}

\smileys

\goodbadpoints

Jaký je širší kontext toho, co učím?\\
\note{Kde v mojí výuce se dotýkám jiné výuky, jiných oblastí vědění?}

\newpage
\marginsection{-90}{1}{volné poznámky}
\vspace*{-2em}
\note{Jak získám od studentů zpětnou vazbu na ty aspekty mé vyúky, které mě zajímají?\\
Jak to dělají ostatní vyučující?\\
A jak jiné předměty?}

% ======= WEEK 13 =======

\chapter{}
\vspace*{-2em}

\marginsection{90}{-1}{před výukou}

Příprava:\hspace{1cm} hodin(y)

Jaký je cíl této hodiny? Proč chci studenty učit tohle?
\vspace*{1cm}

Jak poznám, že se cíle podařilo dosáhnout?

\hspace*{-1cm}
\rule{\rulelength}{0.4pt}

\marginsection{90}{-3}{po výuce}

\smileys

\goodbadpoints

Co mi dal tento semestr výuky?\\
\note{Co jsem si uvědomil(a)? Co jsem zlepšil(a)?}

\newpage
\marginsection{-90}{1}{volné poznámky}
\vspace*{-2em}
\note{Za co jsou studenti v předmětu hodnoceni?\\
Odpovídá to, co zkouška testuje, tomu, co já učím?\\
Pokud mají studenti zájem se dál samostatně rozvíjet v dané oblasti, jsou na to po tomto kurzu připraveni?}

% ======= APPENDICES =======

\restoregeometry
\chapter*{Učitelská rubrika}
\label{rubrika}

Na následujících stranách najdete rubriku činností učitele. Běžně si učitelé-začátečníci (a někdy ani učitelé-veteráni) neuvědomují možné úrovně a dimenze učitelské magie, prostě proto, že se s~nimi nikdy nepotkali. Je možné, že po pár měsících či letech učení přijde moment, kdy si uvědomíte, že máte mnohem větší prostor pro zlepšování, než jste si prve mysleli.\footnote{Více informací najdete například v knize \emph{How Learning Works: Seven Research-Based Principles for Smart Teaching}.}

\section*{Co je to rubrika?}

Rubrika je nástroj pro \emph{self-assessment} (uvědomění a popsání svých vlastních znalostí a dovedností).
Pomáhá zřetelně si uvědomit a pojmenovat, jaké části daná věc má (např. umět programovat znamená být schopen navrhnout algoritmus, efektivně pracovat s~pamětí, znát syntaxi, \dots) a taky na jaké úrovni jsem já v jednotlivých částech (např. navrhnout algoritmus umím dobře, zapsat to kódem mám trochu problém a práci s pamětí neoptimalizuji vůbec).

\section*{Jak učitelskou rubriku používat?}

Na začátku semestru si rubriku vyplň (označ, kde na dané škále jsou podle tebe tvoje schopnosti). Vyber si 1--3 oblasti, na které se chceš zaměřit v tomto semestru a chceš-li být hodně důsledný(á), můžeš vymyslet i konkrétní kroky, které uděláš a indikátory svého postupu. Na konci semestru si můžeš rubriku vyplnit znovu a zhodnotit, v čem jsi se posunul(a).

Rubriku můžeš čist i jako manifest -- popisy kategorie \uv{guru} vystihují naši představu o tom, co zvládají skvělí učitelé a co bychom rádi mezi cvičícími na fakultě kultivovali.

\newcounter{rubricquestion}

\newpage
\rubricpage{Vědomá pozornost a rozhodování, \textit{tracking}}
{Nedokážu si při výuce ani po ní vybavit, co se děje a jaký mám záměr. Při výuce se často cítím ztraceně.}
{Někdy se mi daří uvědomovat si, jaký mám právě záměr, co právě dělám a jaký to bude mít efekt. Většinou ale nemám při výuce vědomou pozornost.}
{Během výuky jsem si téměř vždy vědom(a) toho, co se se mnou i se skupinou děje, co dělám a jaký to bude mít efekt, kam směřuji. Dokážu si vybavit, jak jsme se dostali tam, kde právě jsme.}
{Interakce se skupinou, kladení otázek skupině}
{Se skupinou neinteraguji. Neptám se, asi bych nedostal(a) odpověď. Nevím, jak studenty dobře zapojit.}
{Vím, že lze interagovat se skupinou a znám nástroje (chápu, jak by to šlo). Nedokážu je ale často používat. Někdy  se skupiny ptám, ale nedostávám odpověď.}
{Se skupinou často a efektivně interaguji, ptám se způsobem, který studenty aktivizuje a zapojuje. Situace, kdy nedostávám ze skupiny odpověď dokážu obratně řešit (např. přeformulováním otázky).}

\rubricpage{Strukturování výuky}
{O strukturování výuky nijak nepřemýšlím.}
{Chápu smysl přehledného strukturování své hodiny a snažím se o to. Často se ale zamotám, ztratím nit, nebo řeším mnoho naráz a studenti se pak ztrácí nebo odpojují.}
{Moje hodiny mají jasnou strukturu. Studenti ví, co se právě děje, co bude následovat a chápou návaznosti. Mezi jednotlivými bloky vědomě dělám zřetelné přechody.}
{Vlastní pocit z výuky}
{Svou výuku nijak nereflektuji.}
{Ve výuce si často moc nevěřím, vyžaduje to hodně energie, nebo se cítím napjatě. Mám strach z toho, na co se studenti zeptají.}
{Ve výuce se většinou cítím uvolněně a sebevědomě, mám svůj styl a baví mě to.}

\rubricpage{Formativní zpětná vazba}
{O dávání zpětné vazby studentům nijak nepřemýšlím.}
{Snažím se studentům zpětnou vazbu dávat. Myslím si však, že jí není dost, nebo to nedělám efektivně, nebo to studenti nevnímají jako podporu a projev respektu.}
{Se studenty ve výuce interaguji tak, že dostávají průběžně formativní zpětnou vazbu. Chápou tedy, co jim jde, kde dělají chyby a jak se mohou zlepšovat. Zároveň se moji studenti cítí být respektováni a zpětné vazby se nebojí.}
{Jasné zadávání úloh}
{O zadávání úloh nijak zvlášť nepřemýšlím.}
{Stává se mi, že zadám úlohu a studenti neví, co dělat, jak začít nebo k čemu mají dojít (co má být výsledkem).}
{Když zadávám úlohu nebo aktivitu, mají studenti jasno v tom, co mají dělat a neřeší tak zbytečně věci, které nemám záměr procvičovat.}

\rubricpage{Design výuky a variabilita aktivit}
{Učím tak, jak mi řekli nebo kopíruji výuku, kterou jsem viděl(a). Nepřemýšlím o~jiných variantách.}
{Jsem si vědom(a) toho, že existuje mnoho typů aktivit, které lze při výuce použít. Neznám jich ale dostatek, nedokážu je efektivně zadávat, nebo nemám jasno v tom, proč je vybírat.}
{Znám dostatek typů aktivit, svou výuku skládám tak, aby byla dostatečně pestrá. Vybrané aktivity efektivně učí/procvičují to, co mám záměr učit/procvičovat. Vybrané aktivity zároveň studenty efektivně zapojují a přispívají k jejich motivaci.}
{Širší kontext výuky a konkrétní hodiny}
{O širším kontextu výuky nepřemýšlím.}
{Uvědomuji si, že nemám pro sebe pojmenované znalosti a dovednosti, které u~studentů rozvíjím. Nevím, jak jejich pokrok sledovat. Nevím, v jakém kontextu to studenti využijí.}
{Mám jasnou představu o tom, k čemu studenty vedu, jaké dovednosti rozvíjím, jaké znalosti jim chci předat. Vím, proč tyto dovednosti rozvíjím a v jakém kontextu je studenti v budoucnu použijí. Vím, jak jejich pokrok sledovat.}
\vspace*{-1em}

\rubricpage{Jasné vysvětlování}
{Svoje vysvětlování nijak nereflektuji.}
{Když vysvětluji, běžně se mi stává, že si nejsem jistý, zda vysvětluji dobře a zda to studentům pomáhá při pochopení.}
{Když vysvětluji teorii, demonstruji řešení nebo opravuji chybný postup. Dělám to efektivně a dokážu se dobře vžít do toho, jak to vidí student. Nestává se mi, že by studenti moje vysvětlení nechápali, nebo že bych vysvětloval(a) něco, na co se neptali.}
{Nastavení prostředí, systémy ve výuce}
{O nastavení atmosféry nepřemýšlím, systémy ve své výuce nevnímám.}
{Přemýšlím o nastavení pravidel i atmosféry. Systémy ve výuce (např. bodování) přebírám od ostatních. Nemám ale jasno v efektech, nebo neznám způsoby, jak je lze upravit.}
{Umím ve výuce vytvořit prostředí, které podporuje efektivní učení. U systémů, které používám (např. bodování, bonbóny, zahajovací rituály) chápu efekt. Systémy nepřebírám slepě, chápu jejich účel a uzpůsobuji podle svých potřeb.}

\rubricpage{Improvizace, přizpůsobování}
{V průběhu své výuky na vzniklé situace nijak vědomě nereaguji.}
{Uvědomuji si chvíle, kdy by mohlo být zajímavé nebo užitečné dělat něco jiného, než jsem měl(a) v plánu. Většinou ale nedokážu v daný moment vhodně zareagovat.}
{Dokážu svoji výuku průběžně přizpůsobovat tomu, co se právě děje ve skupině a co studenti potřebují. Mám k tomu dostatek nástrojů a dokážu je efektivně použít.}
{Individuální interakce se studenty}
{O konzultacích nijak vědomě nepřemýšlím.}
{Stává se mi, že si nevím rady při individuální interakci se studentem (např. u~tabule, konzultace). Interakce neprobíhá efektivně nebo se student cítí zastrašen.}
{Při interakcích s jednotlivci (např. u tabule, konzultace) efektivně využívám čas. Studenti se mnou konzultují rádi, cítí ze mě respekt a podporu.}

\rubricpage{Management skupiny, \textit{subgrouping}}
{Přijde mi, že dělit skupinu je zbytečnost, o~práci v~menších skupinkách nepřemýšlím.}
{Uvědomuji si možnosti práce ve dvojicích či menších skupinách. Tuším, že bych toho mohl(a) lépe využívat, hledám jak.}
{Mám jasno v tom, kdy chci pracovat se skupinou jako celkem, kdy s jednotlivci a kdy se skupinkami. Rozdělení do menších skupin efektivně používám. Ve vhodných případech zadávám interakce mezi skupinami.}
{Čtení skupiny}
{Skupinu ve své výuce nesleduji, pozornost věnuji pouze obsahu.}
{Uvědomuji si, že mi skupina vysílá signály a že by bylo dobré jim rozumět a využít je pro efektivní vedení výuky. Ve výuce to ale dokážu jen výjimečně.}
{Dokážu dobře odhadnout naladění skupiny. Mám jasno v tom, co se ve skupině studentů děje (např. únava, rezignace, zájem).}

\chapter*{Indikátory}
\label{indikatory}
\vspace{-0.5em}
\note{Ukazatel stavu; aspekt výuky, který je možné přímo sledovat.}

\section*{Kvantitativní indikátory}

Práce s časem\\
\note{O kolik minut jsem prodlužoval(a)? Kolik času zůstalo navíc?\\
Jak dlouho jsem měl(a) monolog? Kolik času studenti něco dělali?\\
Kolik času jsem v průměru věnoval(a) jednomu studentovi?}

Interakce se studenty\\
\note{Kolik otázek jsem položil(a)? Kolik z nich bylo uzavřených/otevřených?\\
Kolik otázek jsem dostal(a)? Kolik jsem dokázal(a) zodpovědět?\\
Kolik studentů jsem pochválil(a) (za jakkoliv malý úspěch)?}

Řešení úloh\\
\note{Kolik úloh vyřešili studenti na hodině? Kolik úkolů dostali na doma?\\
Kolik studentů bylo ztraceno (těžký úkol)? Kolik se naopak nudilo?}

Účast\\
\note{Kolik studentů přišlo na výuku? Kolik z nich meškalo?\\
Kolik studentů odešlo během výuky?}

\section*{Kvalitativní indikátory}

Pocity a subjektivní dojmy\\
\note{Jaké pocity jsem měl(a) během výuky?\\
Jací byli studenti? Byli aktivní, poslouchali mě, \dots?\\
Dávám svojí prací studentům dobrý příklad?}

Organizace\\
\note{Věděli studenti, co přesně měli dělat, jak, a proč?\\
Do jaké míry jsem dodržel(a) časový plán?}

Obsah výuky\\
\note{Je moje výuka variabilní (různé typy úloh, pomůcek, \dots)?\\
Odpovídá obsah mojí výuky tomu, co chci naučit?}

\newpage

\chapter*{Bloomova taxonomie}
\label{bloom}
\vspace{-0.5em}
\note{Hierarchie kognitivních vzdělávacích cílů, B.\ Bloom, 1956}
\vspace{-0.3em}

\begin{enumerate}[leftmargin=*]
\item \textbf{Vzpomenout si (\textit{remember})}\\
\note{termíny a fakta, jejich klasifikace a kategorizace}\\
{\small definovat, vyjmenovat, zopakovat, popsat, rozpoznat, přiřadit pojem ke konceptu}

\item \textbf{Pochopit (\textit{understand})}\\
\note{překlad do jiné formy, jednoduchá interpretace, extrapolace}\\
{\small interpretovat, reprezentovat, vysvětlit, parafrázovat, sumarizovat, uvést příklad, ilustrovat na příkladu, porovnat}

\item \textbf{Použít (\textit{apply})}\\
\note{užití postupu ve správné situaci, použití abstrakcí a zobecnění}\\
{\small vykonat, aplikovat, využít, demonstrovat v praxi, implementovat, vyřešit uzavřený a ohraničený problém}

\item \textbf{Analyzovat (\textit{analyze})}\\
\note{rozklad na části, vztahy a interakce mezi částmi}\\
{\small diskutovat, \uv{rozbít} na části, rozlišit, navrhnout/zvolit řešení, organizovat, strukturovat, integrovat, propojit}

\item \textbf{Vyhodnotit (\textit{evaluate})}\\
\note{posouzení na základě stanovených kritérií a standardů}\\
{\small zkontrolovat, detekovat, testovat, posoudit, odborně kritizovat, argumentovat}

\item \textbf{Vytvořit (\textit{create})}\\
\note{vytvoření nového celku, reorganizace do nové struktury}\\
{\small generovat, produkovat, vymyslet, plánovat/projektovat, komponovat, designovat, budovat, kreativně kombinovat prvky}
\end{enumerate}

\note{Hranice mezi úrovněmi nejsou striktně určené. Seznam akcí ti ale pomůže popsat cíle hodiny i znalosti a dovednosti, které by měli studenti získat. Ukáže ti také, jestli je tvá výuka dostatečně variabilní.}

\chapter*{Užitečné nástroje}
\vspace{-0.5em}
\note{Připomenutí nástrojů a konceptů z hodin Praktika vedení cvičení}

\section*{Výuka jako celek}

\begin{itemize}
\item Používání jmen -- \note{jmenovky, tahák, ptaní se}
\item Uspořádání prostoru -- \note{rovnost, tvar usazení, pozice učitele}
\item Verbalizace očekávání -- \note{očekávají obě strany to stejné?}
\item Precedenty -- \note{historie, co se stalo a může se opakovat}
\item Zakázky -- \note{co studenti chtějí dostat?}
\item Hospitace -- \note{jak to dělají jiní? Jak to dělám já?}
\item Žádost o zpětnou vazbu -- \note{je to skutečně tak, jak si myslím?}
\item Bloomova taxonomie -- \note{verbalizace cílů hodiny/výuky}
\item Rubrika -- \note{self-assessment, sledování pokroku}
\end{itemize}

\section*{Struktura hodiny}

\begin{itemize}
\item Situování -- \note{shrnutí kde jsme a kam jdeme}
\item Označení přechodů -- \note{explicitní přechod do dalšího bloku}
\item Veřejný checklist -- \note{viditelná struktura hodiny}
\item Check-in -- \note{otevření, zjištění stavu studentů}
\item Check-out -- \note{uzavření, výzva k reflexi}
\item Tracking -- \note{sledování/uvědomování si vývoje diskuze/hodiny}
\item Parkoviště otázek -- \note{odkládání větších otázek na později}
\end{itemize}

\section*{Zadávání úkolů}

\begin{itemize}
\item Otázky ke skupině -- \note{srozumitelnost, podmínka + akce}
\item Multiple choice questions -- \note{uvěřitelnost možností}
\item Hlasování -- \note{pozitivní, negativní, rozdělování bodů, Kahoot}
\item Jasnost zadání -- \note{začátek, proces, výsledek, hodnocení, trvání}
\item Praktická ukázka -- \note{aby se museli i dívat, nejen poslouchat}
\item Rozcvička -- \note{fyzické/myšlenkové probuzení, opakování}
\item Nastavování prahu -- \note{jak moc snadné je to udělat?}
\item Chunking -- \note{dělení na malý počet stravitelných části}
\item Externí motivace -- \note{body navíc, bonbóny, ...}
\end{itemize}

\section*{Tvorba aktivit}

\begin{itemize}
\item Schody -- \note{postupnost malých, ztěžujících se úkolů}
\item Netradiční zadání -- \note{např.\ programování na mnoho způsobů}
\item Argumentace -- \note{zdůvodnění svého názoru, hodnocení jiných}
\item Figurky -- \note{fyzické objekty a manipulace s nimi}
\item Antiproblém -- \note{řešení opačného problému}
\item Myšlenkový oblak -- \note{sbírání asociací, diskuze, souvislosti}
\item Konceptová mapa -- \note{vizualizace pojmů/znalostí a jejich vztahů}
\end{itemize}

\section*{Uvádění aktivit}

\begin{itemize}
\item Subgrouping -- \note{dělení skupiny na menší části}
\item Harvest/sklizeň -- \note{sdílení názorů, zesílení důležitých signálů}
\item Think-pair-share -- \note{nejprve sami, diskuze v páru, sklízení}
\item Peer review -- \note{studenti si navzájem hodnotí práci}
\item Peer tutoring -- \note{studenti se učí navzájem}
\item Neupozornění na chybu -- \note{studenti si ji najdou sami}
\item Obrácená hodina -- \note{samostudium teorie, společné cvičení}
\end{itemize}

\chapter*{Prostor pro poznámky}
\note{Například nejdůležitější zpětná vazba z hospitací.}

\chapter*{Prostor pro poznámky}
\note{Na co chci myslet, když učím?}

\newpage
\vspace*{\fill}
Verze: \today

Autoři deníku: Martin Ukrop, Ondráš Přibyla, Valdemar Švábenský, Martin Macák

Vytvořeno pro kurz \textit{DUCIT Praktikum vedení cvičení} na Fakultě informatiky MU.

\textit{Budeme rádi, když nám dáš vědět svoje postřehy z používaní deníku a další náměty emailem Martinovi na mukrop@mail.muni.cz.}
\end{document}
